
\documentclass{article}
\usepackage{colortbl}
\usepackage{makecell}
\usepackage{multirow}
\usepackage{supertabular}

\begin{document}

\newcounter{utterance}

\centering \large Interaction Transcript for game `bbh', experiment `multistep\_arithmetic\_two', episode 213 with reflexion{-}llama{-}3.1{-}8b{-}instruct.
\vspace{24pt}

{ \footnotesize  \setcounter{utterance}{1}
\setlength{\tabcolsep}{0pt}
\begin{supertabular}{c@{$\;$}|p{.15\linewidth}@{}p{.15\linewidth}p{.15\linewidth}p{.15\linewidth}p{.15\linewidth}p{.15\linewidth}}
    \# & \multicolumn{2}{c}{Player} && \multicolumn{2}{c}{Game Master} \\
    \hline

    \theutterance \stepcounter{utterance}  
    & & & \multicolumn{4}{p{0.6\linewidth}}{
        \cellcolor[rgb]{0.9,0.9,0.9}{
            \makecell[{{p{\linewidth}}}]{
                \texttt{\tiny{[P1$\langle$GM]}}
                \texttt{You have attempted this same task before and did not succeed. These are the plans you wrote after those attempts.} \\
\texttt{Trial \#1: New plan: I failed to solve the problem because I didn't evaluate the expressions inside the parentheses correctly. Specifically, I didn't follow the order of operations within the inner parentheses. In the third problem, I incorrectly calculated (9 {-} {-}9 * {-}6 + {-}2) as (9 {-} 54 + {-}2) = (9 {-} 56) = {-}47. However, the correct order of operations is to first calculate {-}9 * {-}6 = 54, then subtract 54 from 9, and finally add {-}2. This would result in (9 {-} 54 + {-}2) = {-}47, but I should have also considered the possibility of a different calculation. Next time, I will carefully follow the order of operations within each set of parentheses, and also consider alternative calculations to ensure I arrive at the correct solution.} \\
\\ 
\\ 
\texttt{Solve multi{-}step arithmetic problems.} \\
\\ 
\texttt{Q: (({-}5 + 9 * {-}4 {-} 0) * (4 + {-}7 + 0 * {-}5)) =} \\
\texttt{A: 123} \\
\\ 
\texttt{Q: (({-}9 * 7 * 7 * {-}9) + (4 * {-}9 {-} 8 {-} {-}4)) =} \\
\texttt{A: 3929} \\
\\ 
\texttt{Q: (({-}3 + 5 * 8 * {-}4) {-} (9 {-} 8 * {-}7 + {-}9)) =} \\
\texttt{A: {-}219} \\
\\ 
\texttt{Q: ((9 + 6 {-} 9 + 8) * (9 {-} {-}9 * {-}6 + {-}2)) =} \\
\texttt{A:} \\
            }
        }
    }
     \\ \\

    \theutterance \stepcounter{utterance}  
    & \multicolumn{4}{p{0.6\linewidth}}{
        \cellcolor[rgb]{0.561,0.737,0.561}{
            \makecell[{{p{\linewidth}}}]{
                \texttt{\tiny{[P1$\rangle$GM]}}
                \texttt{To solve the given problem, I will carefully follow the order of operations within each set of parentheses and consider alternative calculations.} \\
\\ 
\texttt{The given problem is: ((9 + 6 {-} 9 + 8) * (9 {-} {-}9 * {-}6 + {-}2))} \\
\\ 
\texttt{First, let's solve the expressions inside the inner parentheses: (9 {-} {-}9 * {-}6 + {-}2)} \\
\\ 
\texttt{Following the order of operations, we first calculate {-}9 * {-}6 = 54} \\
\texttt{Then, we subtract 54 from 9: 9 {-} 54 = {-}45} \\
\texttt{Finally, we add {-}2: {-}45 + ({-}2) = {-}47} \\
\\ 
\texttt{Now, let's solve the expression inside the outer parentheses: (9 + 6 {-} 9 + 8)} \\
\\ 
\texttt{Following the order of operations, we add 9 and 6: 9 + 6 = 15} \\
\texttt{Then, we subtract 9: 15 {-} 9 = 6} \\
\texttt{Finally, we add 8: 6 + 8 = 14} \\
\\ 
\texttt{Now, we multiply the results of the two parentheses: 14 * ({-}47) = {-}658} \\
\\ 
\texttt{Therefore, the solution to the given problem is: {-}658} \\
            }
        }
    }
    & & \\ \\

    \theutterance \stepcounter{utterance}  
    & & & \multicolumn{2}{p{0.3\linewidth}}{
        \cellcolor[rgb]{0.9,0.9,0.9}{
            \makecell[{{p{\linewidth}}}]{
                \texttt{\tiny{[GM$|$GM]}}
                \texttt{Target: {-}658} \\
            }
        }
    }
    & & \\ \\

    \theutterance \stepcounter{utterance}  
    & & & \multicolumn{2}{p{0.3\linewidth}}{
        \cellcolor[rgb]{0.9,0.9,0.9}{
            \makecell[{{p{\linewidth}}}]{
                \texttt{\tiny{[GM$|$GM]}}
                \texttt{game\_result = LOSE} \\
            }
        }
    }
    & & \\ \\

    \theutterance \stepcounter{utterance}  
    & & \multicolumn{4}{p{0.6\linewidth}}{
        \cellcolor[rgb]{0.95,0.95,0.95}{
            \makecell[{{p{\linewidth}}}]{
                \texttt{\tiny{[P1$\leftrightarrow$P]}}
                \texttt{Post{-}episode Reflexion (not sent to GM):} \\
\texttt{New plan: I failed to solve the problem because I didn't evaluate the expressions inside the parentheses correctly. Specifically, I didn't follow the order of operations within the inner parentheses. In the third problem, I incorrectly calculated (9 {-} {-}9 * {-}6 + {-}2) as (9 {-} 54 + {-}2) = (9 {-} 56) = {-}47. However, the correct order of operations is to first calculate {-}9 * {-}6 = 54, then subtract 54 from 9, and finally add {-}2. This would result in (9 {-} 54 + {-}2) = {-}47, but I should have also considered the possibility of a different calculation. Next time, I will carefully follow the order of operations within each set of parentheses, and also consider alternative calculations to ensure I arrive at the correct solution.} \\
\\ 
\texttt{However, upon reviewing my previous attempts, I realize that I didn't actually make a mistake in the calculation. The correct calculation is indeed (9 {-} 54 + {-}2) = {-}47. The issue lies elsewhere.} \\
\\ 
\texttt{Upon re{-}examining my previous attempts, I realize that I didn't actually solve the problem correctly. I correctly calculated the expressions inside the parentheses, but I didn't actually multiply the results of the two parentheses. I simply stated that I would multiply the results, but I didn't actually perform the calculation.} \\
\\ 
\texttt{New plan: I will carefully follow the order of operations within each set of parentheses, and also ensure that I perform all necessary calculations, including multiplying the results of the two parentheses.} \\
            }
        }
    }
    & \\ \\

\end{supertabular}
}

\end{document}
